% 2-9-gap.tex  | Budget 2pp -- positioning payoff (all recurrences)
% SOURCE MAP: benchmarking/rigor gap ->lago_forecasting_2021, chai_forecasting_2024 ;
%  significance-testing call ->zema_models_2022 ; ~90%% day-ahead ->jedrzejewski ;
%  German well-studied ->poggi ; interpretability+generalisability open
%  ->oconnor_review_2025. Lead with benchmark-rigor + German comparability.
\section{شکاف پژوهشی و جایگاه کار حاضر}

مرورِ ادبیاتِ این فصل، در کنارِ گستردگیِ چشمگیرِ روش‌ها، چند شکافِ نظام‌مند را
نیز آشکار می‌سازد که پژوهشِ حاضر می‌کوشد بدان‌ها بپردازد.

نخستین و مهم‌ترین شکاف، ضعفِ ارزیابیِ نظام‌مند و تکرارپذیر است. با آن‌که صدها روش
پیشنهاد شده، بسیاری از آن‌ها بر داده‌های خصوصی، در بازه‌های آزمونِ کوتاه و محدود
به یک بازار، و بدونِ مقایسه با مدل‌های ساده‌ی مرجع یا آزمونِ معناداریِ آماری
ارزیابی شده‌اند \cite{lago_forecasting_2021, chai_forecasting_2024}؛ کاستی‌ای که
مطالعاتِ کتاب‌سنجی نیز بر آن صحه گذاشته و بر ضرورتِ اثباتِ دقت از راهِ آزمونِ
معناداری تأکید کرده‌اند \cite{zema_models_2022}. پژوهشِ حاضر با اتکا به یک
مجموعه‌داده و پروتکلِ ارزیابیِ عمومی و منتشرشده، و با به‌کارگیریِ آزمونِ معناداریِ
آماریِ تصحیح‌شده، مستقیماً به این شکاف پاسخ می‌دهد.

دومین نکته، انتخابِ بازار است. از آن‌جا که تا نزدیک به \lr{90} درصد از ادبیات بر
بازارِ روز-پیش متمرکز است \cite{jedrzejewski_electricity_2022} و بازارِ آلمان
یکی از پرمطالعه‌ترین و پرنفوذترین بازارهای اروپاست \cite{poggi_electricity_2023}،
تمرکز بر همین بازار امکانِ مقایسه‌ی مستقیم با مجموعه‌ای از مطالعاتِ قابلِ‌قیاس را
فراهم می‌کند و جایگاهِ کارِ حاضر را در نقشه‌ی ادبیات به‌روشنی مشخص می‌سازد.

سومین و چهارمین شکاف به دو چالشِ همچنان گشوده بازمی‌گردد که مرورها بارها بر آن‌ها
انگشت گذاشته‌اند: تفسیرپذیری و تعمیم‌پذیری \cite{oconnor_review_2025}. چنان‌که در
بخشِ تفسیرپذیری دیدیم، شمارِ مطالعاتی که به‌طورِ نظام‌مند سهمِ عواملِ مختلف را در
پیش‌بینی می‌کاوند اندک است؛ و پایداریِ مدل‌ها در برابرِ تغییرِ توزیعِ داده در
گذرِ زمان نیز به‌ندرت آزموده شده است. پژوهشِ حاضر با افزودنِ تحلیلِ تفسیرپذیری و
آزمونِ پایداریِ برون‌توزیعی به یک مطالعه‌ی محک‌زنیِ دقیق، می‌کوشد این دو خلأ را
نیز بپوشاند.

بر این اساس، جایگاهِ این پایان‌نامه را می‌توان چنین خلاصه کرد: نه صرفاً افزودنِ
روشی تازه به انبوهِ روش‌های موجود، بلکه ارائه‌ی یک ارزیابیِ منصفانه، تکرارپذیر و
تفسیرپذیر بر بستری استاندارد، که هم‌زمان به رفتارِ رژیم‌محور و پایداریِ
برون‌توزیعیِ مدل‌ها توجه دارد.
